% Session 4 of "From Paper to Proof: Formalizing Mathematics in Lean with AI".
% Build: Overleaf with the pdfLaTeX compiler (or pdflatex, twice)
\documentclass[10pt, aspectratio=169]{beamer}
\usepackage{bgsmathlean}
\leanfile{Session4.lean}

\title{From Paper to Proof}
\subtitle{Session 4: Your own theorem}
\author{Michail Karatarakis}
\institute{BGSMath Graduate Course \textperiodcentered{} UPC, Barcelona}
\date{19 November 2026}

\begin{document}

\maketitle

%============================================================================
\begin{frame}[fragile]{The course}
\sessionplan{4}

\vspace{0.6em}
Today: one theorem from start to finish, then some of yours, live.
\end{frame}

%----------------------------------------------------------------------------
\begin{frame}[fragile]{The course in one slide}
\renewcommand{\arraystretch}{1.4}
\begin{tabular}{@{}>{\bfseries}p{0.14\linewidth}>{\raggedright\arraybackslash}p{0.8\linewidth}@{}}
\toprule
Read & Brackets, junk values, coercions, quantifier order. The checklist. (Session 1)\\
Write & Definitions with tests. Statements whose hypotheses can hold, and are all needed. (Session 2)\\
Prove & An agent writes the proof against a pinned statement. (Session 3)\\
Check & No \lc{sorry}, the three standard axioms, the statement unchanged. (Session 3)\\
\bottomrule
\end{tabular}

\vspace{0.8em}
\begin{keypoint}
Lean checks the \textbf{proof}. You check the \textbf{statement}.
\end{keypoint}
\end{frame}

%----------------------------------------------------------------------------
\begin{frame}[fragile]{Today}
\begin{enumerate}
\item \textbf{One theorem, start to finish.} Everything from the course, on one small example.
\item \textbf{Your theorems, live.} Most of the session: two or three of the statements you sent, with the same steps.
\item \textbf{At home:} your own theorem, from \texttt{MyTheorem.lean}.
\end{enumerate}
\end{frame}

%============================================================================
\section{One theorem, start to finish}

\begin{frame}[fragile]{The statement}
\begin{columns}[T]
\begin{column}{0.38\textwidth}
\begin{paper}[Theorem]
Every continuous $f : [0, 1] \to [0, 1]$ has a fixed point.
\end{paper}

\begin{itemize}
\item ``$f : [0,1] \to [0,1]$'' becomes two hypotheses: continuity on $[0,1]$, and \emph{maps into} $[0,1]$.
\item Pinned in a \lc{def}, as in session 3.
\end{itemize}
\end{column}
\begin{column}{0.58\textwidth}
\leanlabel
\lean{fixed_stmt}
\end{column}
\end{columns}
\end{frame}

%----------------------------------------------------------------------------
\begin{frame}[fragile]{Test it before proving it\livedemo}
\lean{fixed_tests}

\vspace{0.3em}
\begin{itemize}
\item The hypotheses can hold together: a constant function.
\item The ``maps into'' hypothesis is needed: without it, $x \mapsto x + 1$ is a counterexample.
\end{itemize}
\end{frame}

%----------------------------------------------------------------------------
\begin{frame}[fragile]{The plan, on paper}
\begin{paper}[Proof]
Let $g(x) = f(x) - x$. Then $g(0) = f(0) \ge 0$ and $g(1) = f(1) - 1 \le 0$. By the intermediate value theorem there is $c \in [0, 1]$ with $g(c) = 0$, that is, $f(c) = c$. \hfill$\square$
\end{paper}

\vspace{0.6em}
\begin{itemize}
\item One step needs another theorem: the intermediate value theorem.
\item State it precisely first, in the form the proof needs. Then look for it.
\end{itemize}
\end{frame}

%----------------------------------------------------------------------------
\begin{frame}[fragile]{A conditional proof}
\begin{columns}[T]
\begin{column}{0.58\textwidth}
\lean{conditional}
\end{column}
\begin{column}{0.38\textwidth}
\begin{itemize}
\item The missing step is a \textbf{hypothesis} of the theorem, not a \lc{sorry} inside it.
\item Anyone reading the statement sees exactly what is assumed.
\item The same shape as the conditional formalization of the unit distance disproof (session 3).
\end{itemize}
\end{column}
\end{columns}
\end{frame}

%----------------------------------------------------------------------------
\begin{frame}[fragile]{Closing the gap}
\lean{close_gap}

\vspace{0.3em}
\begin{itemize}
\item Mathlib has the intermediate value theorem: \lc{intermediate_value_Icc'}.
\item If it did not, \lc{main_of_ivt} would still be an honest result, with its assumption in plain view.
\end{itemize}
\end{frame}

%----------------------------------------------------------------------------
\begin{frame}[fragile]{The audit\livedemo}
\begin{columns}[T]
\begin{column}{0.5\textwidth}
\lean{audit}
\end{column}
\begin{column}{0.46\textwidth}
\begin{itemize}
\item \lc{Statement} is unchanged.
\item No \lc{sorry}; only the three standard axioms.
\item The tests passed before the proof was written.
\end{itemize}
\end{column}
\end{columns}

\vspace{0.6em}
\begin{keypoint}
\normalsize Statement, tests, plan, proof, audit. The same steps work for your theorem.
\end{keypoint}
\end{frame}

%============================================================================
\section{Your theorems, live}

\begin{frame}[fragile]{Your theorems\livedemo}
For each statement you sent:
\begin{enumerate}
\item The paper statement, with every hypothesis, including the implicit ones.
\item The Lean statement, drafted with an AI. Read it aloud, as in session 1.
\item The tests: the hypotheses can hold, and each one is needed.
\item The proof: an agent, against the pinned statement.
\item The audit.
\end{enumerate}

\vspace{0.8em}
\begin{keypoint}
\normalsize The question for each one: does the Lean statement say what the paper says? You decide.
\end{keypoint}
\end{frame}

%============================================================================
\section{At home}

\begin{frame}[fragile]{Your own theorem}
Open \texttt{MyTheorem.lean}. It has one section per step:

\vspace{0.4em}
\renewcommand{\arraystretch}{1.25}
\begin{tabular}{@{}>{\bfseries}p{0.04\linewidth}>{\raggedright\arraybackslash}p{0.27\linewidth}>{\raggedright\arraybackslash}p{0.62\linewidth}@{}}
1 & Paper statement & Every hypothesis, including the implicit ones.\\
2 & Definitions & Only if Mathlib has none. Test each one.\\
3 & Statement & Pinned. The agent may not edit it.\\
4 & Tests & Hypotheses can hold; each one is needed; a case you know.\\
5 & Steps & The plan, as precise statements.\\
6 & Proof & You, an agent, or both.\\
7 & Audit & \lc{#print axioms main}.\\
\end{tabular}

\vspace{0.6em}
Steps 1--4 matter most. A well-tested statement with an open proof is a good outcome.
\end{frame}

%----------------------------------------------------------------------------
\begin{frame}[fragile]{When you are stuck}
\begin{itemize}
\item \textbf{A definition is missing.} Search Mathlib first (Loogle, LeanSearch, the documentation). If it really is missing, write it and test it (session 2).
\item \textbf{A step will not go through.} State it as its own lemma, and give the agent only that lemma.
\item \textbf{The agent keeps failing.} Ask it to try the negation, or to look for a counterexample. The statement may be false.
\item \textbf{The gap is real.} Make it a hypothesis, as in \lc{main_of_ivt}. A conditional result with a precise assumption is progress.
\item \textbf{You are lost.} Read the infoview, then ask on the Lean Zulip.
\end{itemize}
\end{frame}

%============================================================================
\section{After the course}

\begin{frame}[fragile]{Where to go next}
\begin{itemize}
\item \textbf{The Lean Zulip chat}, stream \emph{\#new members}: questions of every level are welcome.
\item \textbf{Mathematics in Lean}, for the parts of Mathlib you did not see here.
\item \textbf{Contributing to Mathlib}: a result in the right generality, with Mathlib's names and style, reviewed by maintainers. Start with something small.
\item \textbf{Larger projects}: Tau Ceti roadmaps, Prove2Me missions, or a blueprint of your own.
\end{itemize}
\end{frame}

%----------------------------------------------------------------------------
\begin{frame}[fragile]{The one idea of the course}
\begin{keypoint}
Lean checks the \textbf{proof}. You check the \textbf{statement}.
\end{keypoint}

\vspace{0.8em}
\begin{itemize}
\item AI now writes a large part of the proofs. That makes the statement the part that needs you.
\item Read it, test it, pin it, audit what comes back.
\end{itemize}

\vspace{1.5em}
\centering
{\Large Thank you.}
\end{frame}

\end{document}
